\documentclass[12pt]{article}
\usepackage[utf8]{inputenc}
\usepackage[T1]{fontenc}
\usepackage[brazil]{babel}
\usepackage[a4paper,margin=2.5cm]{geometry}
\usepackage{enumitem}
\setlist[itemize]{noitemsep,topsep=4pt}
\title{Sprint 3 - S3\_17\\JavaScript Assíncrono e APIs}
\author{Matheus Moreli Durão Silva\\Matrícula 12621CBS217\\Cibersegurança - Bacharelado - Noturno}
\date{2026}
\begin{document}
\maketitle

\section*{Entregável}
Aplicação Web que consome API REST. Este documento registra o objetivo do card, a solução desenvolvida e a forma de conferência. Os arquivos de código e as evidências citadas permanecem na mesma pasta deste documento.

\section*{O que foi estudado e aplicado}
\begin{itemize}
\item Promises
\item async e await
\item Fetch API
\item Tratamento de falhas
\end{itemize}

\section*{Como conferir}
Execute \texttt{python api.py} na pasta S3\_19 e \texttt{python -m http.server 8000} nesta pasta. Abra \texttt{http://127.0.0.1:8000}, clique em \textit{Carregar ativos} e confira a lista obtida de \texttt{GET /ativos}. A evidência real de 13/10/2026 está em \texttt{evidencia\_execucao.md}: um ativo fictício foi mostrado com nome e tipo. A interface trata falhas de rede e de resposta, e usa \texttt{textContent} para exibir dados sem interpretar HTML.

\section*{Conclusão}
O artefato foi organizado para demonstrar tanto o conceito estudado quanto sua aplicação prática. As evidências de execução só são consideradas válidas quando produzidas pela ferramenta correspondente.
\end{document}
